\section{Generación automática de reportes dermatológicos}
\label{sec:marco_reportes}

Que un texto pueda deshacer la salida elegida por la regla de decisión obliga a estudiar el reporte por sí
mismo. Esta sección presenta primero qué se le exige a un reporte dermatológico y después cómo lo produce un
\gls{vlm}, de qué maneras falla el texto generado y qué información adicional puede recibir el generador para
fallar menos.

\subsection{El reporte clínico dermatológico}
\label{sec:marco_rep_tarea}

Un reporte clínico dermatológico es el registro escrito del examen de una lesión. Lo redacta el médico que la
examina, queda en la historia clínica y comunica lo observado a otros profesionales, por ejemplo al
especialista que recibe una derivación. Tiene dos partes. La \textbf{descripción} enumera los hallazgos
visibles, y la \textbf{impresión diagnóstica} los resume en una categoría o, si hay duda, en los diferenciales
de la Sección~\ref{sec:marco_tres_salidas}. La descripción usa un vocabulario estandarizado por consenso, con
términos para el color, el patrón, la distribución y las estructuras de la lesión~\cite{kittler2016standardization},
y algunos de esos términos tienen peso diagnóstico propio. En la lista de 7 puntos, por ejemplo, la red de
pigmento atípica, el velo azul-blanquecino y el patrón vascular atípico son los criterios mayores, los que más
suman a favor de un melanoma~\cite{argenziano1998checklist}.

La generación automática de reportes tiene una historia corta. Su antecedente es la descripción automática de
imágenes (\emph{image captioning}), que produce una oración sobre una imagen natural. El reporte médico es más
largo, sigue una estructura y usa un vocabulario técnico. Su generación con redes neuronales se desarrolló
primero en radiología, desde 2018, con modelos que predicen etiquetas de hallazgos y redactan el reporte
oración por oración, y con arquitecturas que recuerdan patrones de redacción entre reportes
parecidos~\cite{jing2018automatic,chen2020generating}. Con los \glspl{vlm} generativos de la
Sección~\ref{sec:marco_vlm_generativos}, la tarea pasó a resolverse con un único modelo que recibe la imagen y
una instrucción~\cite{hartsock2024vision}.

En todos los casos, el entrenamiento requiere pares de imagen y texto clínico, y esos mismos textos sirven
después de referencia para comparar los reportes generados. En dermatología, estos pares son escasos y
recientes, y su texto tiene orígenes distintos: según el conjunto, lo redactaron dermatólogos, se tomó de libros
de la especialidad o lo generó un modelo y lo revisaron especialistas~\cite{shen2025skincare,zeng2025mmskin,yan2025derm1miccv}.
Ese origen condiciona lo que mide la comparación, porque un texto generado por un modelo, aun revisado, no es el
reporte que un dermatólogo habría escrito al examinar la lesión.

\subsection{Calidad descriptiva y fidelidad diagnóstica}
\label{sec:marco_rep_ejes}

Cada parte del reporte da lugar a un requisito para el texto generado. El primero es describir lo que se ve.
La \textbf{calidad descriptiva} es la correspondencia entre lo que el texto describe y lo que la imagen
muestra, y el vocabulario estandarizado permite medirla, porque fija los términos con que una descripción
nombra cada hallazgo (Sección~\ref{sec:marco_eval_texto}).

El segundo requisito es nombrar el diagnóstico correcto. La \textbf{fidelidad diagnóstica} es la coincidencia
entre el diagnóstico que nombra el texto y el real, y es independiente de la calidad descriptiva. Un texto
puede describir correctamente una red de pigmento atípica y cerrar en un nevo cuando la lesión es un melanoma,
y otro puede nombrar el melanoma correcto apoyándose en un velo azul-blanquecino que la imagen no tiene. Un
reporte es correcto solo si cumple las dos condiciones (Figura~\ref{fig:intro_dos_tareas}).

% Las dos condiciones de un reporte clínico correcto (clasificar y describir una misma imagen).
% Se trajo desde el Capítulo 1 a sec:marco_rep_ejes (2.6.2). El nombre del archivo y la etiqueta
% (fig:intro_dos_tareas) se conservan.

\begin{figure}[H]
    \centering
    \begin{tikzpicture}[x=1mm, y=1mm]

        \node[figentrada, text width=27mm] (img) at (12, 0) {Imagen\\dermatoscópica};

        \node[figtarea, text width=36mm] (t1) at (58, 14)
            {\textbf{Clasificar}\\asignar una categoría};
        \node[figtarea, text width=36mm] (t2) at (58, -14)
            {\textbf{Describir}\\redactar el texto\\de la lesión};

        \node[figsalida, text width=34mm] (cat) at (114, 14)
            {Categoría diagnóstica\\[2pt]{\footnotesize se evalúa con la\\fidelidad diagnóstica}};
        \node[figsalida, text width=34mm] (txt) at (114, -14)
            {Descripción de la lesión\\[2pt]{\footnotesize se evalúa con la\\calidad descriptiva}};

        \draw[figflecha] (img.east) -- ++(6, 0) |- (t1.west);
        \draw[figflecha] (img.east) -- ++(6, 0) |- (t2.west);
        \draw[figflecha] (t1.east) -- (cat.west);
        \draw[figflecha] (t2.east) -- (txt.west);

        \begin{scope}[on background layer]
            \node[figgrupo={las dos condiciones de un reporte correcto},
                  fit=(t1) (t2)] {};
        \end{scope}

    \end{tikzpicture}
    \caption{Los dos requisitos de un reporte clínico correcto: la categoría diagnóstica que nombra se evalúa por su fidelidad diagnóstica, y la descripción de la lesión, por su calidad descriptiva.}
    \label{fig:intro_dos_tareas}
\end{figure}


El tercer requisito viene de la Sección~\ref{sec:marco_incertidumbre}: la certeza del texto tiene que
corresponder a la salida que eligió la regla de decisión. Si la regla afirma un diagnóstico, el reporte lo
sostiene; si ofrece diferenciales, el reporte enuncia los candidatos sin inclinarse por uno que la evidencia no
respalda; y si deriva el caso, el sistema no entrega un diagnóstico. Los tres requisitos dependen de cómo se
produce el texto.

\subsection{Estrategias de generación}
\label{sec:marco_rep_generacion}

Hay dos estrategias para generar un reporte con un \gls{vlm}. En la generación \textbf{de extremo a extremo},
el modelo recibe la imagen $x$ y una instrucción, y produce el texto completo, $y \sim p_\theta(\cdot \mid x)$. Describe la lesión e infiere el diagnóstico a partir de lo que
ve, de modo que un solo modelo tiene que ser a la vez un buen clasificador y un buen redactor.

% Figura: dos formas de generar un reporte con un VLM, de extremo a extremo o condicionado en evidencia
% externa. Se usa en sec:marco_rep_generacion.

\begin{figure}[H]
    \centering
    \begin{tikzpicture}[x=1mm, y=1mm]

        % (a) De extremo a extremo.
        \node[font=\small\bfseries, anchor=west] at (-6, 30) {(a) De extremo a extremo};
        \node[figentrada, text width=22mm] (a_img) at (8, 18) {Imagen e\\instrucción};
        \node[figtarea, text width=26mm] (a_vlm) at (62, 18) {\gls{vlm}};
        \node[figsalida, text width=26mm] (a_rep) at (122, 18) {Reporte};
        \draw[figflecha] (a_img) -- (a_vlm);
        \draw[figflecha] (a_vlm) -- (a_rep);

        % (b) Condicionado en evidencia externa. Arriba, el componente que aporta la evidencia y la regla de
        % decisión; abajo, el mismo recorrido de (a) con la evidencia entrando al VLM desde arriba.
        \node[font=\small\bfseries, anchor=west] at (-6, 4) {(b) Condicionado en evidencia externa};
        \node[figtarea, text width=28mm] (b_ev) at (38, -8) {Clasificador o\\recuperación};
        \node[figtarea, text width=34mm] (b_pol) at (84, -8)
            {Regla de decisión\\[1pt]{\footnotesize afirmar, diferencial\\o derivar}};
        \node[figentrada, text width=22mm] (b_img) at (8, -34) {Imagen};
        \node[figtarea, text width=26mm] (b_vlm) at (84, -34) {\gls{vlm}};
        \node[figsalida, text width=26mm] (b_rep) at (122, -34) {Reporte};
        \draw[figflecha] (b_img.north) |- (b_ev.west);
        \draw[figflecha] (b_ev) -- (b_pol);
        \draw[figflecha] (b_pol) -- node[figetiqueta, right=1mm] {evidencia $c(x)$} (b_vlm);
        \draw[figflecha] (b_img) -- (b_vlm);
        \draw[figflecha] (b_vlm) -- (b_rep);

    \end{tikzpicture}
    \caption{Dos formas de generar un reporte: el \gls{vlm} solo o con la evidencia de
    otro componente.}
    \label{fig:marco_generacion_reportes}
\end{figure}


En la generación \textbf{condicionada en evidencia externa}, que la Figura~\ref{fig:marco_generacion_reportes}
compara con la anterior, otro componente aporta información que se agrega al prompt del generador,
\begin{equation}
    y \sim p_\theta\bigl(\,\cdot \mid x,\ c(x)\bigr),
    \label{eq:marco_rep_condicionada}
\end{equation}
donde $c(x)$ es esa evidencia escrita como texto. Puede ser el diagnóstico que predice un clasificador, los
diagnósticos de casos parecidos recuperados de una base (Sección~\ref{sec:marco_rep_recuperacion}) o los
candidatos que deja la regla de la Ecuación~\ref{eq:marco_politica}. El generador conserva la imagen y redacta
la descripción, pero el diagnóstico le llega como dato. Anclar así el texto en evidencia externa se conoce como
\emph{grounding}, y es el aprendizaje en contexto de la Sección~\ref{sec:marco_llm_grandes}, porque el
generador no cambia sus pesos~\cite{brown2020language}.

La generación condicionada separa los requisitos: la fidelidad diagnóstica pasa a depender de la evidencia, y
la calidad descriptiva sigue a cargo del generador. Su costo es que el reporte hereda los errores de la
evidencia, porque un generador que recibe un diagnóstico equivocado lo afirma con la misma fluidez que uno
correcto. En las dos estrategias, además, la instrucción forma parte de la configuración, ya que puede pedir
que el modelo exponga su razonamiento antes de concluir o que entregue solo el reporte. Un razonamiento
explícito enumera hipótesis que después descarta, de modo que, si se lo lee como parte del reporte, el texto
parece nombrar diagnósticos que no sostiene. Aun separando el razonamiento del reporte, el texto puede fallar de
maneras que una métrica de clasificación no detecta.

\subsection{Fallas del texto generado}
\label{sec:marco_rep_fallas}

Una \textbf{alucinación} es un texto fluido cuyo contenido no está sustentado por la entrada, que en este caso
es la imagen. El fenómeno se describió primero en la generación de lenguaje natural~\cite{ji2023survey} y en la
descripción de imágenes, donde los modelos nombran objetos que la imagen no contiene, en parte por el peso de
los patrones del lenguaje~\cite{rohrbach2018object}, y se documentó después en los
\glspl{vlm}~\cite{bai2024hallucination}. En los \glspl{vlm} médicos, un estudio sistemático encontró que los
modelos ajustados al dominio son más propensos a alucinar que los de propósito general, aun cuando rinden
mejor en las tareas médicas estándar~\cite{gu2024medvh}. En medicina, la alucinación es especialmente costosa,
porque un hallazgo inventado se redacta igual que uno real y puede cambiar una conducta clínica.

La falla contraria se llama aquí \emph{hedging}. El modelo redacta descripciones genéricas, compatibles con
casi cualquier lesión, y evita comprometer un diagnóstico, por ejemplo enumerando varias categorías posibles
sin inclinarse por ninguna. Ese texto no contiene errores visibles, pero tampoco informa nada sobre la lesión
concreta. Expresar incertidumbre no es en sí una falla. Se convierte en \emph{hedging} cuando es
indiscriminada, es decir, cuando no depende de cuán incierto sea el caso. Se ha argumentado que los esquemas de
evaluación que premian acertar y no penalizan adivinar incentivan la alucinación en lugar de la expresión de
incertidumbre~\cite{kalai2025hallucinate}. Ofrecer diferenciales cuando la regla de decisión lo indica
(Sección~\ref{sec:marco_tres_salidas}) es, en cambio, lo opuesto al \emph{hedging}.

\begin{table}[H]
    \centering
    \caption{Fallas de un reporte generado y su consecuencia para el lector, de mayor a menor gravedad.}
    \label{tab:marco_fallas}
    \small
    \begin{tabularx}{\textwidth}{@{}>{\raggedright\arraybackslash}p{3.0cm}
        >{\raggedright\arraybackslash\hspace{0pt}}X >{\raggedright\arraybackslash\hspace{0pt}}X@{}}
        \toprule
        \textbf{Falla} & \textbf{Qué hace el texto} & \textbf{Consecuencia para el lector} \\
        \midrule
        Diagnóstico incorrecto & Afirma una categoría distinta de la real
            & Recibe un error que la fluidez del texto hace parecer confiable \\
        \filasep
        Contradicción & Nombra dos o más diagnósticos incompatibles como conclusión
            & No sabe cuál de ellos sostiene el reporte \\
        \filasep
        \emph{Hedging} & Redacta un texto genérico y evita comprometer un diagnóstico
            & No recibe información y no sabe si el modelo duda o no \\
        \filasep
        Repetición & Repite frases en bucle~\cite{holtzman2020curious}
            & Recibe un texto degenerado, sin información útil \\
        \filasep
        Texto insuficiente & Una o dos frases sin descripción de la lesión
            & No recibe la descripción que el reporte debe contener \\
        \bottomrule
    \end{tabularx}
\end{table}

Además de estas dos, el texto generado puede fallar de otras maneras, y cada reporte se puede asignar a su
falla principal (Tabla~\ref{tab:marco_fallas}). Afirmar un diagnóstico incorrecto es la más grave, porque el
lector recibe un error con la forma de un acierto. La contradicción, que nombra dos diagnósticos incompatibles
en el mismo texto, y el \emph{hedging} dejan la decisión en el lector. La repetición, frecuente con la
decodificación voraz de la Sección~\ref{sec:marco_llm_grandes}, y el texto insuficiente no aportan
información. La alucinación queda fuera de esta clasificación, porque puede acompañar a cualquiera de las
fallas, incluso a un diagnóstico correcto, y no se reconoce por la forma del texto. La taxonomía organiza el
análisis de los reportes generados del Capítulo~\ref{cp:res_gen}.

Las métricas de solapamiento con un texto de referencia (Sección~\ref{sec:marco_eval_texto}) miden un parecido
superficial y pueden no detectar si el reporte transmite los hallazgos clínicos
correctos~\cite{hartsock2024vision}. Por eso la fidelidad diagnóstica y el contenido clínico se miden por
separado. Para reducir estas fallas se puede dar al generador información que la imagen sola no aporta, y la
forma más estudiada de obtenerla es la recuperación.

\subsection{La recuperación como fuente de evidencia}
\label{sec:marco_rep_recuperacion}

\phantomsection\label{sec:marco_rag}%
La \gls{rag}, de 2020, agrega al generador un paso de búsqueda en una base externa de documentos. Tiene tres
partes: la base, un recuperador que selecciona los documentos más relevantes para cada entrada y un generador
que produce la salida condicionado en la entrada y en lo recuperado. En la formulación original, los documentos
son textos, el codificador de consultas del recuperador se ajusta junto con el generador, y el generador produce
una salida con cada documento recuperado y las combina, pesándolas por la relevancia que les asignó el
recuperador. Como el conocimiento queda en la base y no en los pesos, la base se puede cambiar o ampliar sin
reentrenar~\cite{lewis2020retrievalaugmented}. En radiología, la recuperación parte de la imagen y recupera
reportes de casos parecidos, y entrenar el recuperador para que prefiera reportes con los mismos hallazgos
clínicos mejora la corrección de lo generado~\cite{sun2025factaware}.

\phantomsection\label{sec:marco_recuperacion_imagenes}%
Con una base de casos etiquetados, lo recuperable son imágenes. Los $k$ vecinos de la
Ecuación~\ref{eq:marco_vecinos} son los casos más parecidos a la consulta en el espacio de un codificador
visual, y su información, como el diagnóstico o la localización de cada uno, es \textbf{evidencia visual}:
dice qué casos ya vistos se parecen a la consulta. Esa información entra en la
Ecuación~\ref{eq:marco_rep_condicionada} como $c(x) = \psi\bigl(\mathcal{N}_k(x)\bigr)$, donde $\psi$ escribe
como texto lo que se sabe de los vecinos. Un antecedente en descripción de imágenes recupera las descripciones
de imágenes parecidas y las entrega como prompt a un modelo de lenguaje, sin entrenar el
recuperador~\cite{ramos2023smallcap}.

Esta variante difiere de la formulación de \textcite{lewis2020retrievalaugmented} en tres aspectos. El
recuperador puede quedar congelado, sin ajustarse junto con el generador. El generador recibe todo el contexto a
la vez y produce una única salida, en lugar de una por documento. Y lo recuperado son etiquetas y atributos de
imágenes, de la misma modalidad que la consulta, y no documentos de texto. Su límite es el que la
Sección~\ref{sec:marco_senales_confianza} señaló para el consenso de los vecinos: los $k$ vecinos se recuperan
aunque la base no tenga casos parecidos a la consulta, y entonces la evidencia que aportan no describe la
lesión.

\phantomsection\label{sec:marco_integracion}%
La evidencia recuperada y las señales de confianza cumplen, por lo tanto, funciones distintas. La recuperación
aporta qué diagnosticar y atiende la fidelidad diagnóstica. La señal no agrega información al texto, sino que
decide cuánto se puede afirmar con la que hay, y atiende el tercer requisito de la
Sección~\ref{sec:marco_rep_ejes}. Integradas, la evidencia condiciona el texto y la regla de decisión elige si
el sistema afirma un diagnóstico, ofrece diferenciales o deriva el caso. Un mismo componente puede cumplir las
dos funciones: los vecinos recuperados son evidencia para el texto y, a la vez, uno de los predictores cuya
concordancia mide la señal.

Clasificadores, generadores y reglas de decisión producen así salidas de tres tipos: categorías, textos y
decisiones de abstención. Cada una exige sus propias métricas, y ninguna de ellas, medida sobre datos de la
misma fuente que los de entrenamiento, garantiza el desempeño del sistema fuera de esa fuente.
